\documentclass[11pt]{article}
\usepackage[T1]{fontenc}
\usepackage{amsmath,amssymb,amsthm}
\usepackage[margin=1in]{geometry}
\usepackage[hidelinks,pdftitle={Exact Preperiod on the Even Boundary Family {4,6,10K+1}},pdfauthor={Xiaojian Zeng}]{hyperref}

\newtheorem{theorem}{Theorem}
\newtheorem{lemma}[theorem]{Lemma}
\DeclareMathOperator{\mex}{mex}

\title{Exact Preperiod on the Even Boundary Family
  $\{4,6,10K+1\}$}
\author{Xiaojian Zeng\\
  PBC School of Finance, Tsinghua University, 100190\\
  \texttt{xiaojianzeng@foxmail.com}}
\date{}

\begin{document}
\maketitle

\begin{abstract}
For the subtraction game with move set
$S_K=\{4,6,10K+1\}$, $K\geq 1$, we determine the complete
Sprague--Grundy word.  Writing $c=10K+1$, it is a finite concatenation
of five fixed blocks followed by the periodic tail
$\overline{02101}$.  Direct, parameter-independent mex tables verify all
block interiors and seams.  The resulting sequence has least eventual
period $5$ and, under the convention that the period is minimized first,
least preperiod $2c+5=20K+7$.  This treats an even boundary family not
covered by the closest ultimately bipartite odd-parameter result.
\end{abstract}

\section{Statement and context}

A subtraction game with finite move set $S\subset\mathbb Z_{>0}$ is the
normal-play impartial game in which a move subtracts an element of $S$
from a nonnegative heap; see Golomb~\cite{Golomb1966}.  Its
Sprague--Grundy function, in the sense of
Sprague~\cite{Sprague1935} and Grundy~\cite{Grundy1939}, is determined by
\begin{equation}\label{eq:mex}
 G_S(n)=\mex\{G_S(n-s):s\in S,\ s\leq n\},\qquad n\geq0.
\end{equation}
Here $\mex$ denotes the least nonnegative integer absent from the set.

For a sequence $G$, we call $p$ its least eventual period if $p\geq1$ is
least such that $G(n+p)=G(n)$ for every sufficiently large $n$.  After
fixing this least $p$, its least preperiod is the least $q\geq0$ for which
the equality holds for all $n\geq q$.  This is the period-first
convention used by Manabe~\cite{Manabe2026}.

Manabe organizes primitive three-move sets $\{a,b,c\}$ by the angle
$\rho=c\bmod(a+b)$ and proves a sufficient criterion for pure
periodicity.  For $(a,b,c)=(4,6,10K+1)$ one has
$\delta=b-a=2$ and $\rho=1$, which is the boundary angle in the
conjectural positive-preperiod cell of his Table~1 and Conjecture~50
\cite{Manabe2026}.  Zhang's nearby ultimately bipartite family
$\{a,a+2,(2a+2)t+1\}$ requires odd $a\geq3$
\cite[Theorem~6.3]{Zhang2024}, and hence does not include $a=4$.
The result below gives the exact word and both minimality statements in
this even case.  This comparison records the delta from the inspected
sources; it is not a claim of priority.

For a finite word $X=X_0\cdots X_{m-1}$, indices on $X$ below are
zero-based.  Concatenation is written without a symbol, $X^j$ denotes
$j$ repetitions, and $\overline X=XXX\cdots$.

\begin{theorem}\label{thm:main}
For every integer $K\geq1$, set $c=10K+1$ and
$S_K=\{4,6,c\}$.  The Sprague--Grundy sequence of $S_K$ has least
eventual period
\[
 p_K=5
\]
and, for this period, least preperiod
\[
 q_K=2c+5=20K+7.
\]
More precisely, with
\begin{align*}
 A&=0000111122, & B&=2331020010,\\
 D&=1121020010, & F&=112102, & W&=02101,
\end{align*}
the complete SG word is
\begin{equation}\label{eq:word}
 G_{S_K}=A^K\,0\,B\,D^{K-1}\,F\,\overline W.
\end{equation}
\end{theorem}

\section{Uniform block verification}

We prove the exact word by checking recurrence~\eqref{eq:mex} at every
index.  Because the recurrence determines successive terms uniquely,
this is also a uniqueness proof.

Before index $c$, only the moves $4$ and $6$ are available.  Direct
calculation for indices $0,\ldots,9$ gives $A$.  At residues
$0,\ldots,9$ modulo $10$, the two predecessor values at offsets $4$ and
$6$ are
\[
 \texttt{11}\quad\texttt{11}\quad\texttt{21}\quad\texttt{21}
 \quad\texttt{02}\quad\texttt{02}\quad\texttt{00}\quad
 \texttt{00}\quad\texttt{10}\quad\texttt{10}.
\]
Their mex values spell $A$ again.  This propagates through $A^K$ and
also gives the final $0$ at index $c-1=10K$.

At $n=c+t$, $0\leq t<10$, the option reached by subtracting $c$ has
value $A_t$.  Substitution of the already determined short-move values
gives, in move order $(4,6,c)$, the ten triples
\begin{equation}\label{eq:Btable}
 \texttt{110}\ \texttt{210}\ \texttt{210}\ \texttt{020}\
 \texttt{221}\ \texttt{301}\ \texttt{321}\ \texttt{131}\
 \texttt{032}\ \texttt{212}.
\end{equation}
Their mex word is $B$.

Suppose next that a $D$ block is present.  Its residue-$j$ short
predecessors are the circular entries $D_{j-4}$ and $D_{j-6}$, while
the $c$-predecessor is $A_j$ (indices on length-$10$ words are reduced
modulo $10$ here).  The option triples are
\begin{equation}\label{eq:Dtable}
 \texttt{000}\ \texttt{020}\ \texttt{100}\ \texttt{000}\
 \texttt{111}\ \texttt{101}\ \texttt{211}\ \texttt{111}\
 \texttt{022}\ \texttt{212},
\end{equation}
whose mex word is $D$.  The first seam is also covered because
$(B_4,\ldots,B_9)=(D_4,\ldots,D_9)$, precisely the suffix visible to
the moves $4$ and $6$; later seams are $D|D$.  When $K=1$ the exponent
$K-1$ is zero, so no $D$ block occurs.

The word $F$ begins at index $2c-1$.  Whether the preceding length-$10$
block is $B$ (when $K=1$) or $D$ (when $K\geq2$), direct substitution
gives
\begin{equation}\label{eq:Ftable}
 \texttt{000}\ \texttt{022}\ \texttt{103}\ \texttt{003}\
 \texttt{111}\ \texttt{100},
\end{equation}
whose mex word is $F$.  Thus~\eqref{eq:word} is verified through index
$2c+4$.

It remains to check the infinite tail, which begins at
$q=2c+5$.  For $n=q+u$ with $0\leq u<6$, the complete option triples
are
\begin{equation}\label{eq:tailseam}
 \texttt{212}\ \texttt{110}\ \texttt{020}\ \texttt{211}\
 \texttt{000}\ \texttt{221}.
\end{equation}
Their mex values are $\texttt{021010}$, the first six entries of
$\overline W$.  The sixth triple is the same for $K=1$ and $K\geq2$:
its $c$-predecessor is respectively the first entry of $F$ or of $D$,
and both are $1$.

For $u\geq6$, write $r=u\bmod5$ and
\[
 h_u=G_{S_K}(c+5+u),
\]
the value reached from $q+u$ by subtracting $c$.  The two short options
are $W_{r+1}$ and $W_{r-1}$, with indices modulo $5$.  The five possible
local calculations reduce to
\[
\begin{array}{c|c|c|c}
r & \text{short-option pair} & W_r &
 \text{condition for mex to equal }W_r\\ \hline
0&(2,1)&0&h_u\ne0\\
1&(1,0)&2&h_u\ne2\\
2&(0,2)&1&h_u\ne1\\
3&(1,1)&0&h_u\ne0\\
4&(0,0)&1&h_u\ne1
\end{array}
\]
Thus it suffices, in every row, to prove $h_u\ne W_r$.

There are three source regions for $h_u$.  If $K\geq2$, the $D$ region
starts at $u=5$; at its actual alignment
\[
 D=\texttt{1121020010},\qquad WW=\texttt{0210102101},
\]
and the two words differ at every position.  The $F$ region starts at
$u=c-6$, which is divisible by $5$, and
\[
 F=\texttt{112102}
\]
differs positionwise from the aligned word $\texttt{021010}$.  Finally,
for $u\geq c$, the $c$-predecessor is itself in the tail and
\[
 h_u=W_{u-c}=W_{r-1}\ne W_r,
\]
because $c\equiv1\pmod5$ and adjacent cyclic entries of $W$ are
unequal.  The regions meet without a gap.  For $K=1$ the $D$ region is
empty, while the $F$ and tail checks meet the six explicitly checked
seam positions.  This proves~\eqref{eq:word} for every $K\geq1$.

\section{Minimality}

The exact word makes $5$ a period from index $q$.  It cannot start at
an earlier index, since the comparison immediately before $q$ fails:
\begin{equation}\label{eq:boundary}
 G_{S_K}(q-1)=F_5=2,\qquad
 G_{S_K}(q-1+5)=W_4=1.
\end{equation}

We also record the elementary period lemma needed to show that $5$ is
least.

\begin{lemma}\label{lem:gcd}
If an infinite word has eventual periods $r$ and $s$, then it has
eventual period $\gcd(r,s)$.
\end{lemma}

\begin{proof}
Assume $r>s$ and that both period equalities hold beyond $N$.  For
$n\geq N+s$,
\[
 G(n+r-s)=G((n-s)+r)=G(n-s)=G(n),
\]
so $r-s$ is an eventual period after discarding a further finite
prefix.  Iterating the Euclidean algorithm proves the claim.
\end{proof}

If there were an eventual period $p<5$, Lemma~\ref{lem:gcd} applied to
$p$ and $5$ would give eventual period $1$.  This is impossible because
$\overline W$ is not eventually constant: every copy of $W$ contains
both $0$ and $2$.  Hence $5$ is the least eventual period.  With this
period fixed, equation~\eqref{eq:boundary} excludes every preperiod
smaller than $q$, completing the proof of Theorem~\ref{thm:main}.

\section*{Computational and AI assistance}

An exact-integer Python checker independently evaluates every finite
mex table displayed above, every pointwise mismatch used in the tail
argument, and recurrence replays for $1\leq K\leq50$ and also for
$K=100$ and $K=1000$.  These runs are reproducibility checks, not
an extrapolation in $K$; the universal quantifier is established by the
finite block reduction.  This separation between a transparent
mathematical reduction and an exact checker is methodologically akin to
the finite-certificate architecture described by Zeng, Liu, Ratnavelu,
and Huat~\cite{ZengEtAl2026Certificates}, although the mathematical
problems and certificates are different.

OpenAI Codex assisted with computational exploration, exact recurrence
checking, proof drafting, adversarial mathematical review, literature
verification, and preparation of this manuscript.  Language-model
output is not used as mathematical evidence; the proof is the symbolic
argument above, with the exact scripts serving only the stated checking
role.

\bibliographystyle{plain}
\bibliography{references}

\end{document}
